\chapter{绪论}
\label{chap:introduction}

\section{研究背景与意义}

合作行为通常要求个体承担成本，而产生的收益却可能由群体共同分享。在这种条件下，减少投入能够提高个体的当期相对收益，群体有利的合作因而未必会由个体决策自然维持。公共物品博弈把这一矛盾抽象为成员投入、公共收益放大和群体共享之间的关系，为分析个体利益与集体合作的冲突提供了一个可计算模型。

空间公共物品博弈进一步把交互限制在局部邻域。同一个个体会同时参与多个相互重叠的博弈组，合作团簇的形成、团簇边界上的背叛行为以及不同局部区域之间的信息传播由此相互作用。已有研究表明，空间结构以及惩罚、声誉等机制都会改变合作演化条件\cite{brandt2003punishment}。但在给定收益和交互结构之后，个体如何根据经验更新行为同样重要：相同的局部反馈，经由不同学习规则解释后，可能产生完全不同的长期轨迹。

强化学习提供了描述这种持续试错过程的自然工具\cite{sutton2018reinforcement}。本文中的每个智能体维护独立 Q 表，根据自身交互经验估计合作和背叛的动作价值，再通过探索与当前价值判断选择行为。与直接规定个体模仿某个邻居不同，邻居信息可以先进入价值更新，再通过后续决策间接影响行为。由此产生一个核心问题：在自身奖励之外，邻居的哪些信息值得被用作学习参考，以及这些信息应当以多大强度影响动作价值？

本文前期已发表工作从可靠信息条件出发，提出基于邻居相对综合奖励与动作一致性的附加价值修正，并考察影响强度、声誉奖励和感知范围对合作过程的作用\cite{kang2025neighbor}。这一机制把较高表现的邻居作为参考，因此隐含了一个重要条件：接收者观察到的动作和奖励能够代表邻居的真实表现。

然而，实际接收端得到的是报告消息，而不是发送者内部状态本身。若一个发送者持续报告较高奖励，最高奖励规则可能反复选择它；若报告动作也发生失真，附加修正甚至可能被推向与真实行为不同的方向。这类报告仍可落在合法数值范围内，因此不能依赖简单的异常值检查处理。本文据此将研究从“如何利用可靠邻居信息”推进到“信息失真后如何选择参考对象并控制其影响”，并在同一空间社会困境中分析这两个问题之间的连续关系。

这一研究的意义主要在于把社会信息对强化学习的影响落实到具体更新环节。通过区分真实博弈、报告消息和价值修正，可以进一步追问合作改善究竟来自奖励设计、参考对象变化还是更新幅度调整，也能够明确方法在何种信息条件下失效。本文关注的是计算模型中的学习与合作机制，仿真结果用于说明所定义环境中的算法行为，不直接外推为真实人群规律或实际分布式系统的部署结论。

\section{相关研究}
\subsection{空间交互与合作演化}

规则格点是研究局部合作演化的经典环境。Perc等系统讨论了空间公共物品博弈中的重叠群组、策略分布、激励机制以及有限规模效应\cite{perc2017statistical}。规则格点能够固定节点度数和交互几何，使学习规则的作用不被网络结构差异混杂；与此同时，单一规模上的结果不足以说明系统规模变化后的表现，因此本文在新增实验中进一步检查冻结参数在不同网格规模上的迁移。

空间合作研究不仅关注局部结构，也关注激励机制如何改变合作条件。Brandt等在空间公共物品博弈中结合惩罚机会和声誉考察合作演化\cite{brandt2003punishment}。其中的声誉围绕惩罚意愿构造，而本文的声誉根据合作和背叛累积，并进入状态或即时奖励，两者在定义和作用位置上并不相同。

对本文而言，空间结构承担两类功能：一类决定收益如何通过博弈组产生，另一类决定个体能够从哪些邻居获得社会信息。二者需要在模型中明确区分。扩大博弈组会直接改变收益分配，而扩大感知范围可以在保持五人博弈组不变时只增加可参考的信息来源。第一项工作正是在固定一阶博弈交互的情况下改变感知范围，从而单独观察额外社会信息对价值更新的影响。

\subsection{强化学习与邻居经验利用}

在社会困境中使用强化学习，并不意味着合作必然形成。Macy和Flache在两人社会困境中采用带期望水平的Bush--Mosteller随机学习模型，表明反馈被如何解释会显著影响合作能否持续\cite{macy2002learning}。这一结论虽然来自不同学习形式，但说明研究合作时不能只关注收益矩阵，还需要明确个体更新规则。

Q-learning利用即时奖励与后继状态价值更新当前动作估计，其经典收敛结果依赖特定的环境、访问和步长条件\cite{watkins1992q}。在多个个体同时学习时，任一智能体面对的其他策略都在变化，局部状态表示还可能把不同空间情形映射到同一状态。因此，本文把 Q-learning作为个体行为更新机制，通过仿真实验研究群体演化，而不直接把单智能体收敛定理套用于整个系统。

社会信息进入强化学习有多种位置。Zheng等利用一阶邻居中合作与背叛人数的关系构造离散状态，在公共物品博弈及其自愿参与扩展中研究合作演化\cite{zheng2024evolution}。这种方法使邻域结构首先影响状态，再决定相应 Q 表项的学习。其奖励只取个体发起的博弈组，而本文累计五个重叠群组收益，因此两者报告的合作阈值并不能直接横向比较。

当邻域信息以连续统计量出现时，也可以通过函数近似进入价值估计。Tamura和Morita把自身上一动作和其他成员上一轮合作比例作为深度 Q 学习输入，考察公共物品博弈中的群组规模效应\cite{tamura2024analysing}。他们改变的是共同参与博弈的成员数，而本文可以在固定五人博弈组时改变感知范围，因而“交互范围”和“信息范围”在两项研究中承担不同作用。

Li等在超图公共物品博弈中进一步利用其他参与者的行为历史构造强化学习机制，并观察到不同合作区间和反协调空间模式\cite{li2025otherregarding}。这些结果表明，社会信息对合作的影响不仅体现在最终合作比例，也可能体现在策略形成的空间结构与时间过程。本文因此同时观察合作轨迹、空间分布以及价值修正本身，而不是仅用单个最终合作率概括机制。

本文第一项工作采用另一种社会信息入口：邻居经验不进入共享 Q 表，也不替代标准 TD 更新，而是在个体完成自身更新之后，根据邻居相对综合奖励和动作一致性附加修正当前已执行动作的价值\cite{kang2025neighbor}。因此，邻域状态、奖励设计和附加修正分别对应三个不同环节：状态决定当前学习上下文，奖励评价本轮自身经验，附加修正则利用其他个体的相对表现改变动作价值。本文主要研究第三种作用，并进一步分析它在消息不可靠时会出现什么问题。

\subsection{不可靠通信与消息影响控制}

当多智能体系统依赖通信协作时，消息何时发送、哪些内容值得使用以及如何处理不可靠信息都是独立的设计问题。Jiang和Lu提出ATOC，通过注意机制决定通信时机并整合群组内部表示\cite{jiang2018attentional}。本文不重新优化通信拓扑或传输时机，而是假定候选消息已经到达接收端，研究这些动作和奖励报告如何参与持续的价值更新。因此，门控为0表示停止该条消息对附加更新的影响，并不等价于节省消息传输本身的成本。

若要把评分解释成“消息为真”的概率，需要相应的统计生成模型。Mitchell等使用高斯过程描述位置相关观测，比较真实报告与不同虚假来源假设，并据此计算近似后验置信度\cite{mitchell2020gaussian}。本文采用的线性评分并不做这种概率建模，它直接由长期合作目标训练，得分代表在当前任务中作为参考对象的优先级；自身经验残差只是其中一个比较线索，而非校准后的真实性概率。

面向可能被操纵的通信内容，已有工作从不同位置实施防御。R-MACRL将恶意消息生成、消息重构与攻防训练结合，试图恢复可用通信内容\cite{xue2022misspoke}；ADMAC则通过可靠性评估和可分解的决策结构限制可疑消息的作用\cite{yu2024admac}。这些方法分别对应内容修复与影响调节。本文更接近后一思路，但研究对象不是一次性的动作决策，而是邻居报告反复进入 Q 值附加更新后对合作轨迹产生的长期影响。

通信防御也可以直接改变连接结构。Ma等提出DMAC，通过识别和屏蔽关键通信通道，再利用这些扰动重新训练策略，以降低系统对少量通道的依赖\cite{ma2025decentralization}。本文则保持候选邻居和广播方式固定，异常发送者仍然持续提供消息，接收端只能决定参考谁以及让所选消息产生多大价值修正。因此，方法评价重点是信息使用规则，而不是链路重构。

还有研究在明确攻击预算和集成条件下给出通信鲁棒性的条件性证书。Sun等通过消息子集消融与动作投票，分析受扰输出与正常消息支持动作之间的关系\cite{sun2022certifiably}。本文不建立同类证书，而是从价值更新式推导单步修正幅度界和 Q 值范围，再通过有限时间实验评价合作过程。前者描述数值更新性质，后者描述任务行为，两种证据承担不同作用。

在训练组织上，Gupta等比较集中控制、独立策略训练与参数共享，并使用共享策略参数和个体局部观测实现分散执行\cite{gupta2017cooperative}。本文同样共享一套信息使用参数，但每个节点仍维护自己的 Q 表和历史经验。训练阶段以完整群体轨迹评价共享评分与门控，执行阶段各节点只读取自身状态和局部消息。共享的是“如何使用信息”的规则，而不是个体动作价值本身。

总体而言，已有研究分别回答了空间结构如何影响合作、社会信息如何进入强化学习以及不可靠通信如何被处理。本文聚焦于三者交叉处的一个更具体问题：当邻居表现被直接用作价值学习信号时，可靠信息如何改变合作过程；当相同更新路径收到失真消息后，又应怎样选择参考对象并约束其影响。这个问题同时决定了本文的方法输入、威胁模型和必要对照，也限定了最终结论能够覆盖的范围。

\section{研究问题与主要内容}

围绕上述问题，本文包含两项前后衔接的研究。第一项工作研究可靠邻居信息如何参与动作价值更新。在标准 Q-learning 基础上，引入基于相对综合奖励和动作一致性的邻居影响机制，分别分析影响强度、奖励组成和感知范围对合作轨迹的作用，并结合空间分布和策略切换观察不同合作过程。这部分工作已发表，第三章按照原模型与实验来源整理其方法和结果。

第二项工作研究相同信息使用路径在报告失真后如何调整。针对最高奖励参考可能“选错对象”和奖励差可能“放大影响”两个环节，构造基于局部消息与自身经验的参考评分器和连续门控。共享控制参数通过完整学习轨迹上的黑盒搜索获得，测试时保持固定，而每个智能体的 Q 表仍在新环境中从头学习。

第二项工作的评价不仅比较合作率，还围绕方法必要性和作用机制设置多层对照。性能层面使用冻结参数后的留出测试比较完整控制器、固定合作规则及结构消融；机制层面进一步考察参考来源、实际修正幅度和特征依赖；范围层面改变消息类型、协同因子、异常比例、状态、感知范围和网格规模。这样的实验组织用于回答三个连续问题：学习控制是否优于简单规则，收益主要来自哪里，以及这种收益在什么条件下消失。

\section{研究思路与技术路线}

论文首先统一两项工作的基础模型和记号，包括空间博弈收益、状态表示、综合奖励和 Q-learning 时序，并显式区分真实交互量与对外报告消息。在此基础上，第三章分析可靠社会信息如何作为附加价值修正进入学习；第四章保持同一更新接口，在消息层加入受控失真，再把原有 NI 拆分成参考选择与影响幅度两个可独立分析的环节。

第二项工作先在开发环境中训练共享控制参数，并通过独立验证选择冻结模型；随后使用预留环境种子进行最终留出测试，评价方法的主要性能。冻结之后，再通过结构消融、实际幅度匹配、特征移除以及跨消息和环境条件实验解释性能来源并确定适用范围。实现正确性检查作为实验前置条件单独完成，不作为论文主结果展开。

两项工作的证据来源也分别处理。第一项工作在缺少逐次运行记录的部分图中，只解释原图和源稿能够直接支持的现象，不从图像重新构造误差区间。第二项工作保存每次独立运行及对应配置，使均值、运行间分布和配对比较能够回溯到原始输出。这样既保持已发表结果的来源边界，也避免把开发验证与最终留出结果混为同等证据。

\section{论文组织结构}

第一章说明研究背景、相关工作以及两项研究之间的递进关系。第二章统一定义空间公共物品博弈、强化学习、消息观测与评价指标。第三章介绍可靠信息条件下的邻居影响机制及其合作演化结果。第四章针对消息失真构造局部参考选择和连续门控，并分析价值更新方向、幅度界及参数训练方法。第五章依次给出消息敏感性、冻结后的最终留出测试、结构与机制分析、适用范围以及计算开销。第六章总结两项工作的主要认识，明确当前证据能够支持的结论及后续研究方向。